\chapter{RESULTS AND DISCUSSION}
\thispagestyle{empty}

This chapter presents the implemented system, the observed results of live
validation runs, and a discussion of the design decisions, deviations, and
limitations of the implementation.

\clearpage

\section{Implemented System}

The completed platform implements every element of the approved scope. A
teacher can, through the real user interface, create a subject, upload a
syllabus and generate a reviewable chapter--topic structure, upload and
process study materials through the OCR pipeline, review and generate study
content, manage a question bank, assemble a paper that conforms to a pattern,
and publish an assessment. A student can view available assessments, attempt
them with a timer and deadline enforcement, submit, and receive an evaluated
result, and can practise flashcards and questions with instant feedback.
Administrators manage members and their roles.

\section{Observed Results}

Results below are the actual recorded outcomes of live end-to-end runs, taken
from the project validation documents.

\subsection{Live demonstration runs}

The demo and syllabus harnesses completed green through the public tunnel
gateway (e.g.\ the syllabus suite reports 39 checks passing with
\texttt{FAIL=0}). The attempts harness verified that during an active
attempt no student-facing payload contains an answer key or explanation.

\subsection{Material Intelligence pipeline}

A live test of a 40-page uploaded material completed enhancement end to end:
the pipeline produced a version-1 enhancement with 516 sections and 116
segments (56 relevant, 27 uncertain, and 33 irrelevant to the syllabus), of
which 83 segments were mapped, generating 796 mappings (606 topic-level and
190 chapter-level). This run is significant because it reproduced the exact
shape that the original segment-mapping constraint rejected; the constraint
was corrected with a type-aware migration, and the fixed pipeline completed
the identical enhancement.

\subsection{Retry and cancellation}

The cancellation workflow (processing $\rightarrow$ cancelling $\rightarrow$
material returned to the retryable \textsc{queued} state) and the retry
workflow were verified live, including the race case where a concurrent retry
and cancel are correctly serialized with a conflict response.

\section{Discussion}

The design decisions held up in practice. Moving all heavy work to workers
kept the API responsive during long OCR and generation jobs. Validating AI
output on both worker and API boundaries meant that malformed model output
was re-requested or rejected deterministically rather than leaking into the
domain. Deterministic extractors with explicit ambiguity reports made
reviewable pipeline steps feasible. Tenant-scoped queries, enforced in
application code, passed the cross-institute isolation suites.

Two proposal-level notes should be recorded honestly. First, the worker stack
consumes RabbitMQ directly over AMQP with the \texttt{pika} client rather than
through a task-queue framework, while the messaging design remains unchanged.
Second, Redis is declared in the stack but is not exercised by application
code, so no caching performance is claimed.

\section{Limitations}

Quantitative evaluations (latency, throughput, OCR accuracy on a labelled
corpus, user studies) were not performed; this report reports functional and
behavioural evidence only. The checking-system features (online answer-sheet
capture, OMR, on-screen marking) and essay grading remain explicitly out of
scope, as do per-attempt mechanics and practice scoring.